由命题~\ref{ms:prop-2-2}，$A^0$ 半正定，每个实列的散度为零，且满足 $(\partial\Phi)A^0=0$。引理~\ref{ms:lemma-3-1} 中的全纯 Hessian 与梯度恒等式表明，$H(|\Phi|^2)$ 和 $G(g(Q))$ 的像包含于 $(\partial\Phi)^*$ 的四元数张成空间。由于 $(\partial\Phi)^*$ 与 $X$ 无关，平均后的流仍作用于这个张成空间为零。平均保持半正定性及守恒性。乘以 $g(Q)$ 也保持半正定性，并且给实列散度只新增与 $G(g(Q))$ 缩并所得的一项，而这一项为零。特别地，$A_{\rm w}:H(|\Phi|^2)=0$。

由引理~\ref{ms:lemma-4-1}，$H(r)$ 的零空间是一条直线，正交于 $G(r)$ 且包含于 $\ker H(|\Phi|^2)$。余子式矩阵的像在该直线内，矩阵半正定，且时间对角项为 $\det H_{\rm fib}(r)>0$，见引理~\ref{ms:lemma-3-2}。由引理~\ref{ms:lemma-3-3}，其实列均无散度。这证明了矩阵的各项断言与 $\mu>0$，空间块特征值的断言由引理~\ref{ms:lemma-4-1} 得到。

为证密度上界，将引理~\ref{ms:wave-cap-estimates} 应用于环带球冠，得到 $\kappa\asymp h^{-1}$ 和 $\sum_\xi |\widehat f(\xi)|\lesssim h^{-n/2}$。命题~\ref{ms:prop-2-2} 中的权重不超过 $(\kappa/|Y|)^{d/2}\lesssim h^{-d}$，因为 $|Y|\asymp s>a\ge h$，且其极化列向量在这些纤维上有界。故 $\|A^0\|\lesssim h^{-n-2d}$。余子式矩阵至多增加 $Ch^{50n}a^{-2}s^{-2d}$，这是相对于 $\mu$ 的增量，见引理~\ref{ms:lemma-4-1}。因此得到 $\mu^{1/d}\le h^{-30n}$，其中 $h$ 充分小。

将定理~\ref{ms:thm-a-1} 应用于空间变量。由引理~\ref{ms:lemma-4-1}，定义函数的 Hessian 正定；背景矩阵 $B_{\rm fib}$ 半正定，且由 $n\ge2$ 可知其行列式为零；密度 $\lambda\mu$ 严格为正。用固定椭球作参数化，得到光滑的域族。定理给出存在性、唯一性及参数依赖的光滑性，并由其障碍函数给出 $u\le0$。最后，$B_{\rm fib}+\alpha I$ 的特征值为 $\beta+\alpha$ 和 $d$ 重的 $\alpha$，得到所述行列式。
